\chapter{Numerical Accuracy and Optimization Algorithms}

\begin{objectives}
Control floating-point and discretization error, solve nonlinear equations and
optimization problems, and report numerical uncertainty alongside statistical
uncertainty.
\end{objectives}

\section{Error taxonomy}

Numerical work contains:

\begin{itemize}
\item model error from an imperfect representation of the world;
\item statistical error from finite data;
\item discretization error from time, space, or quadrature grids;
\item iteration error from stopping an algorithm;
\item floating-point error from finite precision;
\item implementation error from incorrect code.
\end{itemize}

Increasing Monte Carlo paths reduces sampling error but not model or time-step bias.
An optimizer tolerance smaller than uncertainty in input quotes adds digits, not
information.

\section{Floating-point arithmetic}

Machine numbers have finite relative precision. Subtracting nearly equal quantities
causes catastrophic cancellation. Summing values with different magnitudes loses
small terms; compensated summation helps. Exponentials overflow and underflow;
log-sum-exp stabilizes
\[
\log\sum_i\e^{x_i}
=m+\log\sum_i\e^{x_i-m},\quad m=\max_i x_i.
\]

Comparing floats for exact equality is usually inappropriate. A tolerance should
reflect scale and the numerical problem, not a universal constant.

\section{Root finding}

Bisection on a continuous $f$ with opposite signs at $a,b$ converges robustly and
halves interval each iteration. Newton's method
\[
x_{k+1}=x_k-\frac{f(x_k)}{f'(x_k)}
\]
has local quadratic convergence near a simple root but can diverge or leave the
domain. Brent's method combines bracketing reliability and interpolation speed.

Implied volatility inversion is monotone for ordinary European options with positive
vega. First enforce no-arbitrage price bounds, bracket volatility, and handle near
intrinsic options where vega is tiny.

\section{Gradient-based optimization}

Gradient descent is
\[
x_{k+1}=x_k-\eta_k\nabla f(x_k).
\]
Newton uses Hessian curvature:
\[
x_{k+1}=x_k-\nabla^2f(x_k)^{-1}\nabla f(x_k).
\]
Quasi-Newton methods approximate the
inverse Hessian from gradients. Constrained solvers use active-set, interior-point,
or sequential quadratic programming ideas.

A convex problem has powerful global guarantees; a calibrated stochastic-volatility
objective is often nonconvex. Multiple starts, parameter transformations, gradient
checks, and profile plots are necessary.

\section{Automatic differentiation and checks}

Automatic differentiation applies the chain rule to program operations and can
produce exact derivatives up to floating point. It differentiates the implemented
program, including any mistake. Compare analytic, automatic, finite-difference, and
complex-step derivatives on small cases.

Central finite difference has truncation error $O(h^2)$ but rounding error grows as
$h$ shrinks. A convergence plot across $h$ reveals the useful region.

\section{Linear systems}

Never form a matrix inverse merely to solve $Ax=b$. Use a factorization: Cholesky
for positive definite matrices, QR for least squares, SVD for rank deficiency, and
sparse methods for large structured systems. Report condition numbers or residuals.

\section{Problems}

\begin{exercise}[Core: bisection]
Prove the interval-error bound after $n$ bisection steps and determine iterations
needed for a desired implied-volatility tolerance.
\end{exercise}
\begin{exercise}[Core: cancellation]
Derive a numerically stable formula for the smaller root of a quadratic when
$b^2\gg4ac$.
\end{exercise}
\begin{exercise}[Advanced: gradient audit]
Design a derivative test for a constrained portfolio objective and distinguish
absolute from relative tolerance near zero.
\end{exercise}


\chapter{Monte Carlo and Variance Reduction}

\begin{objectives}
Simulate stochastic models, quantify standard errors and time-step bias, use variance
reduction, and price path-dependent claims with reproducible diagnostics.
\end{objectives}

\section{Plain Monte Carlo}

For $V=\E[g(X)]$, iid samples give
\[
\widehat V_N=\frac1N\sum_{i=1}^Ng(X_i),\qquad
\widehat{\mathrm{se}}=
\frac{s_g}{\sqrt N}.
\]
The central limit theorem yields an approximate interval
$\widehat V_N\pm z_{1-\alpha/2}\widehat{\mathrm{se}}$. Halving standard error
requires roughly four times as many independent paths.

Pseudorandom seeds make a run repeatable. Separate random streams prevent small code
changes from silently changing every comparison. Common random numbers improve
precision of differences between scenarios.

\section{Simulating diffusions}

Euler--Maruyama for
$\dd X_t=b(X_t)\dd t+a(X_t)\dd W_t$ is
\[
X_{n+1}=X_n+b(X_n)\Delta t+a(X_n)\sqrt{\Delta t}Z_{n+1}.
\]
Under regularity, strong convergence order is typically $1/2$ and weak order $1$.
Milstein adds derivative terms and can achieve strong order $1$ in scalar settings.

For geometric Brownian motion use its exact transition. Naive Euler can produce
negative prices. For CIR variance, full truncation, quadratic-exponential, or exact
transition schemes handle boundaries more reliably.

\begin{definition}[Strong and weak error]
Strong error measures pathwise approximation such as
$\E[\abs{X_T-\widehat X_T}]$. Weak error measures expectation error
$\abs{\E[g(X_T)]-\E[g(\widehat X_T)]}$.
\end{definition}

Pricing a European payoff needs weak accuracy; hedging, barriers, and exposure paths
can require stronger path fidelity.

\section{Antithetic variables}

If $Z$ is standard normal, simulate $g(Z)$ and $g(-Z)$ and average. Variance is
\[
\frac12\Var(g(Z))+\frac12\Cov(g(Z),g(-Z))
\]
relative to a pair-based convention. Negative covariance reduces error. Antithetics
are especially effective for monotone payoffs under symmetric shocks.

\section{Control variates}

Let $Y$ have known mean $\mu_Y$. Use
\[
\widehat V_c=\overline X-c(\overline Y-\mu_Y).
\]
The variance-minimizing coefficient is
$c^\star=\Cov(X,Y)/\Var(Y)$ and reduction factor is $1-\rho_{XY}^2$.
Estimate $c$ on an independent pilot or account for its estimation.

For an arithmetic Asian call, a geometric Asian or discounted terminal stock is a
natural control.

\section{Importance sampling}

If rare loss $A$ is seldom observed under $p$, sample from $q$ and weight:
\[
\Prob_p(A)=\E_q\left[
\ind_A(X)\frac{p(X)}{q(X)}
\right].
\]
The support of $q$ must cover relevant $p$ states. Poor likelihood ratios can have
infinite variance even when the estimator is unbiased. Exponential tilting chooses
a proposal that makes the rare event less rare.

\section{Quasi-Monte Carlo}

Low-discrepancy sequences fill the unit cube more uniformly than random points.
Randomized Sobol sequences allow error estimation. Effective dimension matters:
Brownian bridge or principal-component construction assigns early coordinates to
important variation.

Classical $N^{-1}$-like convergence requires smoothness and dimension conditions;
discontinuous barriers and exercise rules reduce gains. Scrambling and replicate
randomizations provide honest empirical errors.

\section{Nested simulation}

CVA, capital, and risk of future value may require outer market scenarios and inner
conditional valuations. Cost grows multiplicatively. Regression proxies, least-squares
Monte Carlo, multilevel Monte Carlo, and allocation of inner paths reduce cost.
Proxy validation must focus on tails and boundary states, not only average error.

\section{Problems}

\begin{exercise}[Core: control variate]
Derive $c^\star$ and the variance-reduction factor.
\end{exercise}
\begin{exercise}[Core: discretization study]
Design a grid refinement study that separately estimates Monte Carlo error and weak
time-step bias.
\end{exercise}
\begin{exercise}[Advanced: importance sampling]
For a Gaussian tail probability, derive exponential mean shifting and the likelihood
ratio. Explore variance as the shift changes.
\end{exercise}


\chapter{Finite Differences, Trees, and Transform Methods}

\begin{objectives}
Discretize pricing PDEs, analyze stability and convergence, use Fourier methods, and
choose an algorithm that respects payoff and model structure.
\end{objectives}

\section{Finite-difference grids}

For Black--Scholes PDE, transform or discretize time $t_n$ and price $S_j$. Central
spatial differences approximate
\[
V_S\approx\frac{V_{j+1}-V_{j-1}}{2\Delta S},\qquad
V_{SS}\approx\frac{V_{j+1}-2V_j+V_{j-1}}{\Delta S^2}.
\]
Explicit time stepping computes the next layer directly but has a stability
restriction. Implicit stepping solves a tridiagonal system and is unconditionally
stable in a norm sense. Crank--Nicolson averages explicit and implicit operators,
giving second-order time accuracy for smooth solutions.

Nonsmooth terminal payoffs can make Crank--Nicolson oscillate. Rannacher smoothing
uses initial implicit half steps. Grid boundaries should follow asymptotic payoff
behavior and be stress-tested.

\section{Consistency, stability, convergence}

\begin{theorem}[Lax equivalence, informal]
For a well-posed linear initial-value problem and a consistent finite-difference
scheme, stability is equivalent to convergence.
\end{theorem}

Consistency means local truncation error vanishes. Stability controls amplification
of perturbations. Convergence means numerical solution approaches the true solution.
Checking only a visually smooth surface proves none of them.

\section{American complementarity}

An American option solves a linear complementarity problem:
\[
\max\{h-V,\ V_t+\mathcal LV-rV\}=0
\]
under a sign convention. Projected SOR, penalty methods, and policy iteration enforce
the obstacle. Verify value exceeds intrinsic, exercise boundary behaves sensibly,
and European limits are recovered.

\section{Lattices}

Binomial and trinomial lattices discretize state transitions while preserving
recombination. Trinomial trees can match drift and variance with greater flexibility.
For multiple factors, tree size grows quickly and correlation matching can create
negative transition probabilities.

Trees naturally handle early exercise and discrete dividends. They converge
nonmonotonically for some strikes; Richardson extrapolation and smoothing help.

\section{Fourier pricing}

If a model has known characteristic function of log price, option values can be
computed by Fourier inversion. Damping makes a call transform integrable. The
Carr--Madan style transform applies FFT across a strike grid:
\[
C(k)=\frac{\e^{-\alpha k}}{\pi}
\int_0^\infty
\operatorname{Re}\left[
\e^{-iuk}\Psi(u;\alpha)
\right]\dd u.
\]

FFT is fast for many evenly spaced log strikes but couples strike spacing and
frequency cutoff. Aliasing, truncation, damping, and interpolation must be checked.
Direct quadrature can be preferable for a few options.

\section{Numerical quadrature}

Trapezoidal and Simpson rules approximate regular finite integrals. Gaussian
quadrature chooses nodes and weights exact for polynomial families under a weight.
Adaptive quadrature allocates evaluations where estimated error is high.

Singular integrands, oscillation, and infinite domains require transformations or
specialized schemes. Report absolute error against a closed form when one exists.

\section{Algorithm selection}

\begin{longtable}{p{0.22\textwidth}p{0.31\textwidth}p{0.37\textwidth}}
\toprule
Method & Strength & Main limitation\\
\midrule
Closed form & speed, transparency, Greeks & narrow assumptions\\
Tree & early exercise, discrete events & dimension and convergence\\
Finite difference & PDEs, free boundaries & state dimension, grids\\
Monte Carlo & high dimension, paths & early exercise, noisy Greeks\\
Fourier & many strikes, affine models & transform and grid errors\\
Quadrature & low-dimensional accuracy & dimension and singularities\\
\bottomrule
\end{longtable}

A production library prices benchmarks by two independent methods. Agreement within
tolerance is a control, not redundant work.

\section{Problems}

\begin{exercise}[Core: explicit stability]
Derive coefficients of an explicit Black--Scholes scheme and a sufficient condition
for nonnegative weights.
\end{exercise}
\begin{exercise}[Core: boundary]
Specify lower and upper boundaries for European calls and puts with dividends.
\end{exercise}
\begin{exercise}[Advanced: FFT]
Derive the damped call transform and analyze how log-strike spacing relates to the
frequency grid.
\end{exercise}


\chapter{Statistical Computing and Reproducible Software}

\begin{objectives}
Design auditable quantitative software, test mathematical properties, preserve data
lineage, and make research reproducible across machines and time.
\end{objectives}

\section{Architecture}

Separate:

\begin{itemize}
\item immutable source-data ingestion;
\item validation and canonical schemas;
\item pure mathematical functions;
\item stateful simulation or backtest engines;
\item experiment configuration;
\item reports and generated artifacts.
\end{itemize}

Pure functions are easy to test. Domain objects should carry units, currency,
calendar, and convention rather than relying on hidden globals. A pricing function
should not silently download data.

\section{Testing layers}

Unit tests cover small functions and edge cases. Property tests verify invariants:
call price increases with spot, discount factors are positive, a covariance estimate
is positive semidefinite, and portfolio weights reconcile. Regression tests freeze
trusted examples. Integration tests exercise complete data-to-report flows.

Numerical tests use tolerances tied to a benchmark and grid. Tests should include
zero time, zero volatility limits, deep moneyness, negative rates when supported,
missing dates, and invalid inputs.

\section{Data contracts}

A canonical dataset records:

\begin{itemize}
\item provider and exact endpoint;
\item instrument identifier and mapping history;
\item event, publication, and ingestion timestamps with timezone;
\item units, currency, frequency, adjustment, and missing-value codes;
\item download time, query, checksum, and license note;
\item transformations from raw to analytical columns.
\end{itemize}

Never overwrite raw files. A transformation manifest makes derived data reproducible.
CSV is portable but loses types; Parquet preserves schemas and is efficient. Database
constraints can enforce uniqueness and temporal validity.

\section{Environments and randomness}

Pin direct and transitive dependencies for a released result. Record interpreter,
operating system, BLAS, and compiler where numerical differences matter. Containers
improve portability but do not guarantee permanent access to external data.

A random seed initializes a generator. Parallel simulations need independent
substreams, not repeated identical seeds. Store seed and generator family in
experiment metadata.

\section{Logging and experiment manifests}

An experiment manifest includes configuration, Git commit, input checksums, start
time, code version, environment, and output checksums. Logs should report decisions
and anomalies without leaking credentials or licensed data.

Generated tables should come directly from saved results. Manual transcription
breaks lineage. LaTeX can input CSV-derived tables or generated fragments, but the
generation step must be deterministic and documented.

\section{Code review for quantitative models}

Review mathematical specification separately from software behavior. Ask:

\begin{enumerate}
\item Do equations, conventions, and units match the design?
\item Are array shapes, labels, and time axes explicit?
\item Are invalid domains rejected rather than silently coerced?
\item Are calibration residuals and solver status checked?
\item Are data and model timestamps point-in-time?
\item Are independent benchmarks and sensitivities present?
\item Are performance optimizations behavior-preserving?
\end{enumerate}

\section{Problems}

\begin{exercise}[Core: property tests]
Specify property-based tests for Black--Scholes call price, historical Expected
Shortfall, and a covariance shrinkage estimator.
\end{exercise}
\begin{exercise}[Core: manifest]
Design a JSON experiment manifest for a yield-curve study with raw-source checksums
and point-in-time settings.
\end{exercise}
\begin{exercise}[Advanced: reproducibility]
Explain why an unpinned vendor endpoint can make a computationally deterministic
script scientifically irreproducible.
\end{exercise}


\chapter{Modern Machine Learning for Quantitative Finance}

\begin{objectives}
Formulate machine learning as statistical decision theory, understand kernels and
neural networks, build probabilistically calibrated sequence models, use
reinforcement learning cautiously, and evaluate adaptive systems under financial
nonstationarity.
\end{objectives}

\section{Learning as risk minimization}

Let $Z=(X,Y)\sim P$ and prediction rule $f$ lie in class $\mathcal F$. Population
risk is
\[
\mathcal R(f)=\E_P[\ell(Y,f(X))].
\]
Empirical risk minimization chooses
\[
\widehat f\in\argmin_{f\in\mathcal F}
\left\{\frac1n\sum_{i=1}^n\ell(Y_i,f(X_i))
+\lambda\mathcal J(f)\right\}.
\]
The loss defines the statistical target. Squared loss targets conditional mean;
absolute loss targets median; pinball loss targets a conditional quantile; log loss
targets a conditional distribution. A trading objective may instead depend on
cost-aware utility, but directly optimizing a noisy backtest makes selection bias
especially severe.

Expected generalization error decomposes conceptually into irreducible noise,
approximation error, estimation error, and distribution shift. Increasing model
capacity can reduce approximation error and increase estimation error. Regularization,
data scale, and stable inductive structure govern the balance.

\section{Features and invariances}

Prices in currency units are usually nonstationary and arbitrary under stock splits.
Returns, spreads, ratios, ranks, and curve coordinates often encode more stable
objects, but no transformation guarantees stationarity. A feature definition must
include lookback window, missingness rule, publication lag, cross-sectional universe,
and training-only normalization.

Useful inductive biases include:

\begin{itemize}
\item permutation invariance across an unordered asset set;
\item translation invariance across time in local convolutions;
\item monotonicity where economic dominance requires it;
\item no-arbitrage constraints on price surfaces;
\item equivariance to currency or notional scaling;
\item sparsity or low rank for high-dimensional panels.
\end{itemize}

Encoding a known invariance reduces sample burden and prevents economically
impossible extrapolation.

\section{Kernel methods}

A positive definite kernel $k(x,x')$ represents an inner product
$\langle\phi(x),\phi(x')\rangle$ in a possibly infinite-dimensional feature space.
Kernel ridge regression solves
\[
\min_{f\in\mathcal H_k}
\sum_{i=1}^n(y_i-f(x_i))^2+\lambda\norm{f}_{\mathcal H_k}^2.
\]

\begin{theorem}[Representer theorem]
For a broad class of losses depending on training evaluations and a strictly
increasing penalty in RKHS norm, a minimizer has form
\[
\widehat f(x)=\sum_{i=1}^n\alpha_i k(x_i,x).
\]
\end{theorem}
\begin{proof}[Proof idea]
Decompose any $f$ into a component in the span of
$\{k(x_i,\cdot)\}_{i=1}^n$ and an orthogonal component. The latter is zero at every
training point by the reproducing property, so it cannot improve fit and can only
increase the norm penalty.
\end{proof}

For squared loss,
$\widehat\alpha=(K+\lambda I)^{-1}y$. Kernel methods are strong for structured
small-to-medium data but scale cubically without approximation. Kernel and
hyperparameter selection must occur within time-respecting validation.

\section{Neural networks}

A feed-forward network composes affine maps and nonlinear activations:
\[
h^{(0)}=x,\qquad
h^{(\ell)}=\sigma(W_\ell h^{(\ell-1)}+b_\ell),\qquad
\widehat y=W_Lh^{(L-1)}+b_L.
\]
Backpropagation applies reverse-mode automatic differentiation to compute the
gradient of loss with respect to every parameter. Stochastic gradient methods update
using minibatches; momentum and adaptive scaling change optimization dynamics.

Universal approximation results state that sufficiently wide networks can
approximate broad classes of continuous functions on compact sets. They do not say
finite financial samples identify the correct function, that gradient descent finds
it, or that it extrapolates across regimes.

Weight decay, dropout, early stopping, augmentation, and architectural constraints
regularize differently. Batch normalization can leak cross-sectional or temporal
information if a batch mixes future observations; inference statistics and panel
construction require explicit design.

\section{Sequence models}

A recurrent state model updates
$h_t=g_\theta(h_{t-1},x_t)$ and emits a forecast. Gated recurrent units address
vanishing gradients. Temporal convolutions process fixed receptive fields in
parallel.

Self-attention maps queries, keys, and values:
\[
\operatorname{Attention}(Q,K,V)=
\operatorname{softmax}\left(\frac{QK^\transpose}{\sqrt d}\right)V.
\]
Causal masks prevent a position from attending to future tokens. Positional or time
encodings distinguish order and irregular spacing.

Financial panels add complications: changing universes, mixed frequencies, release
lags, missing-not-at-random data, and corporate events. A transformer with a causal
mask can still leak if input features were standardized globally or fundamentals
were aligned to fiscal period rather than publication.

\section{Cross-sectional learning}

At each date, a model can rank assets by a score $s_{i,t}$. Cross-sectional losses
include pairwise ranking and listwise objectives. Portfolio construction converts
scores into weights under neutrality, risk, turnover, and capacity:
\[
w_t=\argmax_w
\left\{
s_t^\transpose w-\frac{\gamma}{2}w^\transpose\widehat\Sigma_tw
-C(w-w_{t-1})
\right\}.
\]
This separation permits evaluation of forecast quality and portfolio translation.
End-to-end optimization can be useful, but it entwines model learning with an
estimated risk and cost layer.

Panel validation blocks by time, not individual rows. Otherwise the model sees the
same market shock in training assets and validation assets. A geography or asset-class
holdout tests cross-sectional transport but does not replace a future-time test.

\section{Probabilistic forecasts and calibration}

A probabilistic model should be sharp and calibrated. For event probability
$p_t$, calibration means
$\Prob(Y_t=1\mid p_t=p)=p$ in an ideal population sense. Reliability diagrams and
Brier decomposition diagnose calibration; log loss heavily penalizes confident
errors.

Quantile forecasts should attain coverage and respond to state. A collection of
independently estimated quantiles can cross; monotone parameterizations or
rearrangement fix ordering. Distributional networks can output location, scale,
mixture weights, or a normalizing flow.

Conformal prediction constructs finite-sample marginal coverage under exchangeability.
Time dependence and distribution shift break the basic guarantee; weighted,
blockwise, or online variants add assumptions rather than magically restoring it.

\section{Uncertainty}

Aleatoric uncertainty describes outcome randomness conditional on inputs. Epistemic
uncertainty describes limited knowledge of the predictive rule. Ensembles, Bayesian
neural approximations, bootstrap, and last-layer methods estimate parts of epistemic
uncertainty, but common model misspecification can make every ensemble member agree
and be wrong.

For a portfolio decision, propagate predictive scenarios through constraints and
costs. Confidence in a score is not the same as confidence in net P\&L because
covariance, impact, and crowding are additional uncertain layers.

\section{Online learning and drift}

Online learning receives data sequentially, predicts, observes loss, and updates.
Regret relative to comparator $f$ is
\[
\operatorname{Regret}_T(f)=
\sum_{t=1}^T\ell_t(f_t)-\sum_{t=1}^T\ell_t(f).
\]
Sublinear regret means average loss approaches the best fixed comparator in hindsight
under the theorem's assumptions. Markets may change so a dynamic comparator or
variation budget is more relevant.

Drift monitoring includes feature distributions, residuals, calibration, turnover,
and economic exposures. A detected change triggers a predeclared action: investigate,
reduce risk, retrain, or retire. Repeatedly changing the system after losses without
an audit creates an adaptive overfit.

\section{Reinforcement learning}

A Markov decision process has state $s_t$, action $a_t$, transition law, reward
$r_t$, and discount $\gamma$. For policy $\pi$,
\[
V^\pi(s)=
\E^\pi\left[
\sum_{k=0}^{\infty}\gamma^kr_{t+k}\mid s_t=s
\right].
\]
The Bellman optimality equation is
\[
V^\star(s)=
\max_a\E[r(s,a,S')+\gamma V^\star(S')\mid s,a].
\]

In trading, the true state is partially observed, actions affect costs and perhaps
prices, rewards have long horizons, and historical data were generated by another
policy. Offline RL cannot freely test counterfactual actions absent support.
Importance sampling can have explosive variance; model-based simulation inherits
simulator error.

A safe design begins with a transparent stochastic-control benchmark, hard risk
constraints, conservative off-policy evaluation, and stress outside the training
distribution. Reward must include costs, financing, inventory, and risk; a raw P\&L
reward invites leverage.

\section{Generative models and synthetic scenarios}

Variational autoencoders, adversarial networks, diffusion models, and autoregressive
models can generate return paths or order-book states. Fidelity should be evaluated
on marginal tails, serial dependence, cross-asset dependence, conditional dynamics,
and decision-relevant stress.

Synthetic data are useful for privacy, augmentation, and controlled experiments.
They are not observed evidence and cannot validate the assumptions used to generate
them. Rare crises are precisely where flexible generators have least information.

\section{Evaluation protocol}

A modern ML study should report:

\begin{enumerate}
\item a simple linear and economic benchmark;
\item model capacity, all tuned ranges, and number of research trials;
\item time and cross-sectional split with purge and embargo if needed;
\item training curves and seed variability;
\item proper statistical loss and calibration;
\item net decision performance with costs and constraints;
\item factor, regime, and subgroup attribution;
\item perturbation, missingness, delay, and cost stress;
\item inference-time latency, stability, and monitoring plan.
\end{enumerate}

Complexity is justified by stable incremental evidence, not by the sophistication of
the architecture.

\section{Problems}

\begin{exercise}[Core: representer theorem]
Complete the orthogonal-decomposition proof of the representer theorem and derive
kernel-ridge coefficients.
\end{exercise}
\begin{exercise}[Core: attention leakage]
List every way a causally masked sequence model can still use future financial
information through preprocessing, data alignment, and universe construction.
\end{exercise}
\begin{exercise}[Advanced: probabilistic portfolio]
Design a model that forecasts a joint return distribution and a policy that optimizes
Expected Shortfall subject to turnover. Specify a proper forecast score and a
separate decision score.
\end{exercise}
\begin{exercise}[Advanced: offline control]
Formulate inventory liquidation as an MDP. State what historical-policy support is
needed for off-policy evaluation and how a model error could create a fictitious
profitable action.
\end{exercise}
\begin{exercise}[Research]
Compare ridge, gradient-boosted trees, and a small causal transformer on one
predeclared target. Use identical point-in-time features, nested walk-forward tuning,
multiple seeds, cost-aware portfolio mapping, and a locked final period.
\end{exercise}
